\section{Related Work}

CrossFlow is related to diffusion and flow-based generation, latent diffusion and autoencoding, and generative models that allow the input and output spaces to differ. The distinction from most prior work is architectural as well as objective-level: CrossFlow keeps the generative input in a compact latent space but supervises the model directly in pixel space.

\subsection{Diffusion, Flow Matching, and Few-Step Generation}

Diffusion models learn to reverse a noise-corruption process. DDPMs~\citep{DDPM} and score-based models~\citep{NCSM} established this paradigm, while DDIM~\citep{DDIM} introduced deterministic sampling for faster inference. Flow matching~\citep{FlowMatching,rf} instead learns a velocity field that transports noise to data along a prescribed path. Transformer architectures such as DiT~\citep{dit} and SiT~\citep{sit} improved scalability and have become common backbones for high-quality image generation. Many methods reduce the number of sampling steps. Consistency models~\citep{CM} learn mappings that remain consistent along probability-flow ODE trajectories, enabling one-step or few-step generation, and improved consistency training~\citep{ICT} stabilizes this objective. Progressive distillation~\citep{ProgressiveDistillation}, distribution matching distillation~\citep{DMD}, MeanFlow~\citep{meanflow}, and improved MeanFlow~\citep{imf} also compress or reformulate iterative sampling into fewer evaluations. CrossFlow shares the goal of efficient generation, but changes the supervised space. These methods typically define the trajectory, velocity, and prediction target in the smae space; CrossFlow takes a noisy latent-space input while predicting a pixel-space image.

\subsection{Latent Diffusion and Autoencoders}

Latent diffusion models~\citep{ldm} reduce generation cost by running diffusion in the latent space of a pretrained autoencoder rather than directly in pixel space~\citep{adm}. This strategy underlies many modern systems, including SDXL~\citep{sdxl}, Stable Diffusion 3~\citep{sd3}, FLUX~\citep{flux}, and SANA~\citep{sana}. The autoencoder is therefore a core part of the system, and methods such as VQ-VAE~\citep{vqvae}, VQ-GAN~\citep{vqgan}, and DC-AE~\citep{dcae} study the trade-off among compression, reconstruction fidelity, and efficiency. A related line of work improves compatibility between autoencoders and latent generators. LSGM~\citep{LSGM} proposed score-based generative modeling in latent space with joint optimization; REPA-E~\citep{REPA-E} tuned VAEs with latent diffusion Transformers through representation alignment; DSD~\citep{DSD} interpreted diffusion as self-distillation for end-to-end latent diffusion; and unified-latent VAEs~\citep{ULVAE,UNITEVAE} co-optimize the tokenizer and diffusion model. These methods address the mismatch between tokenization and generation, but they generally retain separate encoder, generator, and decoder modules. CrossFlow explores a complementary decomposition: the encoder may remain fixed, while the generator and decoder are merged into one model trained on noisy latent inputs and pixel-space targets.

\subsection{Generation Across Different Spaces}

Natural images have strong spatial and semantic structure, so their intrinsic dimension can be lower than their ambient pixel dimension. GANs map low-dimensional noise directly to images~\citep{GAN,biggan,stylegan}, but standard normalizing flows~\citep{NormalizingFlows,RealNVP} require invertible transformations and therefore impose dimension matching. Augmented normalizing flows~\citep{AugmentedNF} and injective flows~\citep{InjectiveFlows} relax this restriction in different ways. Standard diffusion and flow-matching models~\citep{DDIM,DDPM,NCSM,FlowMatching} usually avoid the issue by defining the full trajectory inside a single space. CrossFlow targets the latent-to-pixel case. It does not require an invertible map and does not decode a generated latent through a separately trained module. Instead, the latent path is used only to define noisy inputs and time derivatives, while the supervised prediction is an image. % Cross-modal generation is another setting where input and output spaces differ~\citep{multimodalflow1,multimodalflow2}; many such systems introduce a shared representation between modalities. CrossFlow addresses a narrower but practically important question: how to obtain a flow-style training objective when the noisy input remains latent-valued and the desired output is pixel-valued.
